\chapter{Marchés de l'électricité~I~: conception}\label{ch:m3:power-markets-design}

Le dimanche 11 mai 2025, entre 13~h et 14~h, les panneaux solaires allemands ont produit 44
gigawatts alors que le pays entier en consommait 40. La veille, l'enchère du jour J-1 avait fixé le prix de cette
heure à moins 250 euros le mégawattheure~: les producteurs étaient payés pour livrer
de l'électricité, et les consommateurs étaient rémunérés pour la soutirer. Cinq mois plus tôt, à 17~h, un
jour sans vent, le 12 décembre 2024, la même enchère avait abouti à 936 euros. L'électricité est le seul
produit de base qui doit être produit au moment même où il est consommé, et les marchés de l'électricité sont
conçus autour de ce fait. Ce chapitre décrit cette conception~: l'\omterm{def:m3:power-markets-design:mtu}{unité de temps de marché},
l'enchère du jour J-1 et la manière dont son algorithme détermine le prix, l'\omterm{def:m3:power-markets-design:merit}{ordre de préséance
économique} qui fixe le prix, les zones et les nœuds, ainsi que les \omterm{def:m3:power-markets-design:capacity}{marchés de capacité} et de réserves que les prix de l'énergie
seuls ne rémunèrent pas.

\section{L'électricité en tant que produit de base}

L'électricité se stocke difficilement à grande échelle, et l'offre et la demande sur un réseau doivent
s'égaler en tout instant, faute de quoi sa fréquence s'écarte de sa valeur nominale de 50 ou 60 hertz. Les
marchés échangent donc de l'énergie pour livraison sur des intervalles courts et fixes, et l'opérateur
de réseau équilibre en temps réel ce qui subsiste.

\begin{definition}[Unité de temps de marché]\label{def:m3:power-markets-design:mtu}
L'\emph{unité de temps de marché}\index{unité de temps de marché} (market time unit, MTU) est l'intervalle de
livraison des produits d'un marché de l'électricité~: un produit d'énergie est une puissance constante livrée sur une
MTU, soit une quantité d'énergie exprimée en MWh par MTU.
\end{definition}

Un marché de l'électricité est une suite de marchés portant sur la même MTU~: les forwards et les
contrats à terme pour les mois et les années à venir, l'enchère du jour J-1, le négoce infrajournalier en continu
(\cref{ch:m3:power-markets-trading}), et l'équilibrage en temps réel assuré par l'opérateur. L'enchère du
jour J-1 sert de point d'ancrage~: la majeure partie des volumes en Europe y est valorisée, et la plupart
des forwards se règlent sur ses prix.

\section{L'enchère du jour J-1 et son algorithme de détermination du prix}

\begin{definition}[Marché J-1]\label{def:m3:power-markets-design:da}
Le \emph{marché J-1}\index{marché J-1} (day-ahead) est une enchère journalière à offres scellées
tenue la veille de la livraison, qui fixe pour chaque MTU du lendemain un
prix d'équilibre à prix uniforme (livre~2, chapitre~4) par \omterm{def:m3:power-markets-design:zone}{zone de dépôt des offres}, à partir des ordres
d'achat et de vente soumis avant son heure de fermeture du guichet.
\end{definition}

Pour chaque MTU, les acheteurs soumettent une courbe de demande en escalier (quelle quantité ils achèteront à chaque
prix) et les vendeurs une courbe d'offre en escalier. L'algorithme retient les ordres acceptés
qui maximisent le bien-être, c'est-à-dire la valeur des offres d'achat acceptées diminuée du coût des
offres de vente acceptées, sous réserve des contraintes de réseau entre zones, et fixe dans chaque zone le
prix auquel l'offre et la demande se croisent. Chaque ordre accepté est payé ou paie à
ce prix~; personne ne paie son prix offert.

\begin{proposition}[Le prix d'équilibre soutient l'allocation]\label{prop:m3:power-markets-design:support}
Au prix d'équilibre $p$ d'une zone unique à courbes en escalier, chaque offre d'achat acceptée a une limite
$\ge p$ et chaque offre de vente acceptée une limite $\le p$~; aucune offre d'achat rejetée n'a de limite
supérieure à $p$ et aucune offre de vente rejetée n'a de limite inférieure à $p$. Le volume accepté maximise le bien-être.
\end{proposition}

\begin{proof}
Servir les offres d'achat par ordre de prix décroissant et les offres de vente par ordre croissant, et s'arrêter lorsque
l'offre d'achat suivante vaut moins que l'offre de vente suivante ajoute, chaque unité apportant un surplus positif et aucune
unité un surplus négatif~: c'est le maximum de bien-être. L'ordre marginal, partiellement
exécuté, fixe $p$, qui se situe par construction entre les limites des ordres acceptés et celles des
ordres rejetés de chaque côté.
\end{proof}

Les enchères réelles ajoutent des ordres qui couvrent plusieurs MTU.

\begin{definition}[Ordre bloc]\label{def:m3:power-markets-design:block}
Un \emph{ordre bloc}\index{ordre bloc} (block order) est un ordre d'achat ou de vente d'une quantité dans chacune de
plusieurs MTU, à exécuter en totalité ou pas du tout, assorti d'une limite sur le prix moyen de ces MTU.
\end{definition}

Les blocs permettent à une centrale thermique, avec son coût de démarrage et sa durée minimale de fonctionnement,
d'exprimer son économie réelle dans ses offres. Ils font de la détermination du prix un problème combinatoire~:
un bloc peut être rejeté alors que les prix qu'il aurait obtenus ont une moyenne supérieure à sa limite (un bloc rejeté de
façon paradoxale), mais il n'est jamais accepté lorsque leur moyenne lui est inférieure. Les enchères du jour J-1
européennes sont déterminées par un seul algorithme à l'échelle de tout le continent.

\begin{definition}[Couplage des marchés]\label{def:m3:power-markets-design:coupling}
Le \emph{couplage des marchés}\index{couplage des marchés} (market coupling) est la détermination conjointe des
enchères du jour J-1 de plusieurs zones par un seul algorithme, qui alloue implicitement la capacité
de transport transfrontalière~: l'électricité circule de la zone à bas prix vers la zone
à haut prix jusqu'à ce que les prix convergent ou que la capacité soit épuisée.
\end{definition}

\begin{dated}{2026-09}{Couplage européen du jour J-1}\label{dat:m3:power-markets-design:sdac}
Le Single Day-Ahead Coupling organise une seule enchère à 12~h CET pour une livraison le lendemain,
déterminée par l'algorithme EUPHEMIA dans toutes les \omterm{def:m3:power-markets-design:zone}{zones de dépôt des offres} participantes. Il est passé
d'\omterm{def:m3:power-markets-design:mtu}{unités de temps de marché} de 60 minutes à 15 minutes le jour de négoce du 30 septembre 2025, pour une livraison
à partir du 1\up{er} octobre 2025. Les limites harmonisées du prix d'équilibre étaient de $-\qty{500}{EUR/MWh}$ et
$+\qty{4000}{EUR/MWh}$~; la décision 02/2026 de l'ACER a révisé la méthodologie, et le minimum
a été abaissé à $-\qty{600}{EUR/MWh}$ à partir du jour de négoce du 28 mai 2026.
\end{dated}

\section{L'ordre de préséance économique, les spark spreads et les dark spreads}

\begin{definition}[Ordre de préséance économique]\label{def:m3:power-markets-design:merit}
L'\emph{ordre de préséance économique}\index{ordre de préséance économique} (merit order) est le classement des capacités de production par
coût marginal de court terme, du plus faible au plus élevé~; lorsque les offres sont fixées au coût marginal, la
courbe d'offre de l'enchère est l'ordre de préséance économique et le prix est le coût marginal de la dernière
centrale nécessaire.
\end{definition}

Le coût marginal d'une centrale thermique est son coût de combustible par MWh d'électricité, auquel s'ajoute le coût
des quotas d'émission qu'elle doit restituer (\cref{ch:m3:emissions-and-environmental-markets})~:
$(P_{\mathrm{fuel}} + e\,P_{\mathrm{CO_2}})/\eta$, où $\eta$ désigne son rendement et $e$ les
émissions par MWh de combustible. L'éolien et le solaire n'en ont presque aucune.

\begin{definition}[Spark spread, dark spread]\label{def:m3:power-markets-design:spark}
Le \emph{spark spread}\index{spark spread} d'une centrale à gaz de rendement $\eta$ est le prix de l'électricité
diminué du prix du gaz qu'elle brûle par MWh d'électricité,
soit $P_{\mathrm{power}} - P_{\mathrm{gas}}/\eta$~; le \emph{dark spread}\index{dark spread} est
le même pour une centrale à charbon.
\end{definition}

\begin{omfigure}
\begin{tikzpicture}
\begin{axis}[width=11cm,height=5.2cm,
  xlabel={capacité cumulée, GW},ylabel={coût marginal, EUR/MWh},
  tick label style={font=\small},label style={font=\small},
  xmin=0,xmax=80,ymin=0,ymax=290,grid=major,grid style={gray!25},
  table/col sep=comma]
\addplot[omDef,very thick,const plot] table[x=gw,y=mc]{figdata/markets-3/05-power-markets-design/merit.csv};
\draw[omThm,thick,dashed] (axis cs:45,0) -- (axis cs:45,255);
\draw[omThm,thick,dashed] (axis cs:60,0) -- (axis cs:60,255);
\node[omThm,font=\footnotesize,anchor=east,align=right] at (axis cs:44.2,218) {demande\\45 GW};
\node[omThm,font=\footnotesize,anchor=west,align=left] at (axis cs:60.8,218) {demande\\60 GW};
\node[font=\footnotesize,anchor=south west] at (axis cs:0.3,4) {éolien, solaire};
\node[font=\footnotesize,anchor=south] at (axis cs:25.5,86) {lignite};
\node[font=\footnotesize,anchor=south] at (axis cs:38.5,93) {CCGT};
\node[font=\footnotesize,anchor=south] at (axis cs:56.5,93) {charbon};
\node[font=\footnotesize,anchor=south] at (axis cs:69,132) {OCGT};
\node[font=\footnotesize,anchor=south] at (axis cs:74.5,252) {pétrole};
\end{axis}
\end{tikzpicture}

\omcaption{Un \omterm{def:m3:power-markets-design:merit}{ordre de préséance économique} (merit order) illustratif avec 10~GW d'éolien et de solaire~: gaz à
\qty{35}{EUR/MWh} de combustible, charbon à \qty{12}{}, lignite à \qty{5}{}, carbone à
\qty{70}{EUR/t}. À 45~GW de demande, une centrale à cycle combiné au gaz fixe le prix
(\qty{89.35}{EUR/MWh})~; à 60~GW, une centrale à charbon le fait (\qty{89.68}{}). Données~: le
tutoriel du chapitre.}\label{fig:m3:power-markets-design:merit}
\end{omfigure}

Avec les combustibles de la \cref{fig:m3:power-markets-design:merit} et l'électricité à
\qty{100}{EUR/MWh}, une centrale à gaz de rendement 55~\% gagne un \omterm{def:m3:power-markets-design:spark}{spark spread} de
\qty{36.36}{EUR/MWh} et une centrale à charbon de rendement 40~\%, un \omterm{def:m3:power-markets-design:spark}{dark spread} de \qty{70.00}{} avant
prise en compte du carbone~: après celle-ci, la centrale à charbon est plus chère (\cref{ch:m3:emissions-and-environmental-markets}).
L'ajout d'énergies renouvelables décale tout l'empilement vers la droite et fait baisser le prix à chaque heure
où elles produisent~: avec 30~GW d'éolien et de solaire et 60~GW de demande, le prix passe
de \qty{89.68}{} à \qty{89.35}{EUR/MWh}~; avec 55~GW, à \qty{10}{EUR/MWh}.

\section{Tarification zonale et nodale}

Un réseau a des limites, et un prix doit indiquer où il s'applique.

\begin{definition}[Zone de dépôt des offres]\label{def:m3:power-markets-design:zone}
Une \emph{zone de dépôt des offres}\index{zone de dépôt des offres} (bidding zone) est la plus grande zone géographique à l'intérieur de laquelle
les participants au marché peuvent échanger de l'énergie sans allouer de capacité de transport~; le
\omterm{def:m3:power-markets-design:da}{marché J-1} (day-ahead) fixe un seul prix par zone et par \omterm{def:m3:power-markets-design:mtu}{unité de temps de marché} (MTU).
\end{definition}

L'Europe fixe ses prix par zones, pour l'essentiel nationales~; les marchés organisés des États-Unis fixent des prix par
nœuds.

\begin{definition}[Prix marginal local, rente de congestion]\label{def:m3:power-markets-design:lmp}
Le \emph{prix marginal local}\index{prix marginal local} (LMP) en un nœud du
réseau est le coût de la fourniture d'un MWh supplémentaire en ce point, au moindre coût, sous réserve des
limites du réseau~: sur les marchés américains, la somme d'un prix de l'énergie du système,
d'une composante de congestion et d'une composante de pertes. La \emph{rente de congestion}\index{rente de congestion} sur une ligne contrainte est
l'écart de prix entre ses extrémités multiplié par le flux qui la traverse~; elle revient à quiconque
détient les droits de capacité de la ligne (dans les marchés couplés, aux opérateurs de réseau).
\end{definition}

\begin{proposition}[Deux zones, une ligne]\label{prop:m3:power-markets-design:twozones}
Deux zones compensées conjointement par une ligne de capacité $C$ ont des prix égaux si le flux
qui les égalise ne dépasse pas $C$. Sinon, le flux vaut $C$, de la zone bon marché vers la
zone chère, chaque zone est compensée isolément avec ce flux, qui joue le rôle d'un export ou d'un import pour un preneur de prix, et
la \omterm{def:m3:power-markets-design:lmp}{rente de congestion} vaut $(p_B - p_A)\,C$.
\end{proposition}

\begin{proof}
La compensation conjointe maximise le bien-être total. Si le flux de l'optimum sans contrainte respecte la
capacité de la ligne, il est réalisable et optimal, avec un seul prix. Sinon, le bien-être est concave en fonction du flux, donc
l'optimum contraint fixe le flux à sa borne~; chaque zone est alors compensée avec ce flux fixé,
et les prix diffèrent du prix d'ombre de la ligne.
\end{proof}

\begin{omfigure}
\centering
\begin{tikzpicture}[font=\small,
  zone/.style={draw,thick,rounded corners=3pt,align=center,minimum width=3.6cm,minimum height=1.8cm},
  arr/.style={-{Stealth[length=2.5mm]},very thick}]
\node[zone,draw=omDef,fill=omDef!8] (a) at (0,0) {zone A~: venteuse\\40 GW de renouvelables\\demande 50 GW};
\node[zone,draw=omThm,fill=omThm!8] (b) at (8,0) {zone B~: calme\\5 GW de renouvelables\\demande 55 GW};
\draw[arr,omProp] (a.east) -- node[above=3pt,font=\footnotesize] {ligne saturée~: 3 GW} (b.west);
\node[omProp,font=\footnotesize,anchor=north,align=center] at (4,-0.3) {rente de congestion~:\\\qty{7.30}{EUR/MWh} $\times$ 3 GW};
\node[font=\footnotesize,anchor=north] at (0,-1.15) {prix \qty{82.38}{EUR/MWh} (lignite)};
\node[font=\footnotesize,anchor=north] at (8,-1.15) {prix \qty{89.68}{EUR/MWh} (charbon)};
\node[font=\footnotesize,anchor=north] at (4,-1.85) {avec une ligne de 20~GW~: un seul prix, \qty{89.35}{EUR/MWh} (CCGT)};
\end{tikzpicture}

\omcaption{Deux zones couplées, avec l'\omterm{def:m3:power-markets-design:merit}{ordre de préséance économique} de
la \cref{fig:m3:power-markets-design:merit}. Une ligne de 3~GW est saturée et les prix se séparent~; les
opérateurs de réseau encaissent environ \money{EUR}{21900} pour l'heure. Avec une ligne de 20~GW, les prix convergent. Données~: le
tutoriel du chapitre.}\label{fig:m3:power-markets-design:zones}
\end{omfigure}

La tarification zonale masque la congestion à l'intérieur d'une zone, que l'opérateur doit alors résorber
par un redispatch, en rémunérant certaines centrales pour produire moins et d'autres pour produire plus~; la tarification nodale la révèle
dans le prix, mais elle donne des milliers de prix au lieu d'un seul, et expose chaque producteur à
la congestion locale de son nœud. Le chapitre suivant montre comment les traders couvrent cette exposition.

\section{Marchés de capacité, services système et prix négatifs}

Un marché de l'énergie seule ne rémunère les producteurs que pour les MWh qu'ils vendent. Une centrale qui ne tourne que quelques
heures par an doit amortir ses coûts fixes sur ces seules heures, à des prix de rareté que les régulateurs
plafonnent souvent.

\begin{definition}[Marché de capacité, service système]\label{def:m3:power-markets-design:capacity}
Un \emph{marché de capacité}\index{marché de capacité} (capacity market) rémunère des ressources pour leur disponibilité à
produire (ou à réduire leur consommation) pendant les futures périodes de tension sur le système, au moyen d'enchères
tenues plusieurs années à l'avance, avec des pénalités en cas de défaillance. Un \emph{service
système}\index{service système} (ancillary service) est un service autre que l'énergie que l'opérateur de réseau
achète pour assurer la sécurité du système~: des réserves de fréquence qui réagissent en quelques secondes à quelques minutes,
un soutien de tension, et la capacité à relancer le réseau après une panne généralisée.
\end{definition}

Le plafond de prix et le \omterm{def:m3:power-markets-design:capacity}{marché de capacité} sont deux réponses au même problème. Le marché texan,
qui ne repose que sur l'énergie, a fixé son prix à son plafond de \qty{9000}{\$/MWh} pendant plusieurs jours de février
2021, lors d'une tempête hivernale~; l'ordonnance du régulateur a été contestée devant les tribunaux et confirmée par la
Cour suprême du Texas en 2024. Le surveillant du marché a établi que le maintien des prix au plafond pendant 32
heures après la fin des délestages, jusqu'au matin du 19 février, a ajouté 16~milliards de dollars américains
aux coûts du marché.

À l'autre extrémité, les prix deviennent négatifs lorsque l'offre qui ne peut pas s'arrêter, ou qui ne veut pas le faire, excède
la demande augmentée de la capacité d'exportation. Trois types de producteurs continuent de produire à prix négatifs~:
les centrales dont l'arrêt et le redémarrage coûtent plus que la perte liée au fonctionnement pendant quelques heures~;
les producteurs qui reçoivent une subvention par MWh produit et qui produisent tant que le prix augmenté de la
subvention reste positif~; et les centrales qui doivent fonctionner pour d'autres raisons (chaleur pour une ville,
sécurité du système). Le 11~mai 2025, la production solaire allemande a dépassé la consommation du pays pendant cinq
heures (\cref{fig:m3:power-markets-design:may11}). Dans les données allemandes, 457~heures ont eu un prix moyen négatif
en 2024 et 573 en 2025 (moyennes horaires, telles que les comptabilise aussi
la Bundesnetzagentur~; à partir d'octobre 2025, l'enchère se règle au quart d'heure), presque toutes au
milieu de la journée (\cref{fig:m3:power-markets-design:neghours}).

\begin{omfigure}
\begin{tikzpicture}
\begin{axis}[width=11cm,height=4.8cm,axis y line*=left,
  xlabel={heure, heure locale allemande},ylabel={EUR/MWh},
  tick label style={font=\small},label style={font=\small},
  xmin=-0.5,xmax=23.5,ymin=-300,ymax=200,grid=major,grid style={gray!25},
  xtick={0,3,6,9,12,15,18,21},
  table/col sep=comma]
\addplot[omThm,very thick,const plot mark mid] table[x=hour,y=price]{figdata/markets-3/05-power-markets-design/may11.csv};\label{plot:m3:may11price}
\end{axis}
\begin{axis}[width=11cm,height=4.8cm,axis y line*=right,axis x line=none,
  ylabel={GW},ylabel near ticks,
  tick label style={font=\small},label style={font=\small},
  xmin=-0.5,xmax=23.5,ymin=0,ymax=75,
  legend columns=3,legend style={font=\footnotesize,at={(0.5,-0.33)},anchor=north,draw=none,fill=none,/tikz/every even column/.append style={column sep=8pt}},
  table/col sep=comma]
\addlegendimage{/pgfplots/refstyle=plot:m3:may11price}\addlegendentry{prix day-ahead}
\addplot[omProp,thick] table[x=hour,y=solar]{figdata/markets-3/05-power-markets-design/may11.csv};\addlegendentry{production solaire}
\addplot[omDef,thick,dashed] table[x=hour,y=load]{figdata/markets-3/05-power-markets-design/may11.csv};\addlegendentry{consommation}
\end{axis}
\end{tikzpicture}

\omcaption{Allemagne, dimanche 11~mai 2025, valeurs horaires~: le prix day-ahead (axe de gauche) est tombé à
$-\qty{250.32}{EUR/MWh}$ à 13~h, tandis que la production solaire (axe de droite) a dépassé la consommation de 11~h
à 16~h. Données~: Bundesnetzagentur | SMARD.de, CC BY 4.0.}\label{fig:m3:power-markets-design:may11}
\end{omfigure}

\begin{omfigure}
\begin{tikzpicture}
\begin{axis}[width=11cm,height=4.4cm,ybar,bar width=7pt,
  xlabel={heure, heure locale allemande},ylabel={heures négatives},
  tick label style={font=\small},label style={font=\small},
  xmin=-0.7,xmax=23.7,ymin=0,ymax=190,grid=major,grid style={gray!25},
  xtick={0,3,6,9,12,15,18,21},
  table/col sep=comma]
\addplot[fill=omDef!60,draw=omDef] table[x=hour,y=count]{figdata/markets-3/05-power-markets-design/neghours.csv};
\end{axis}
\end{tikzpicture}

\omcaption{Heures day-ahead allemandes à prix négatif, 2024 et 2025 réunies (1\,030
heures), par heure locale de livraison~: les heures solaires dominent. Données~: Bundesnetzagentur |
SMARD.de, CC BY 4.0.}\label{fig:m3:power-markets-design:neghours}
\end{omfigure}

\section{Tutoriel~: calage d'une heure et couplage de deux zones}\label{tut:m3:power-markets-design:clear}

\begin{tutorial}
\textbf{Objectif.} Caler une MTU à partir d'ordres par paliers, puis coupler deux zones par une ligne.
\textbf{État final~:} les prix des \cref{fig:m3:power-markets-design:merit,fig:m3:power-markets-design:zones}
et les heures négatives de la \cref{fig:m3:power-markets-design:neghours}.
\begin{tutsteps}
\item \textbf{Une zone.} Appariement en \omterm{def:m3:power-markets-design:merit}{ordre de préséance économique} et prix de l'ordre marginal.
\omcode{code/firm/dayahead/firm_dayahead.py}{30}{58}{Calage à prix uniforme d'une unité de temps de marché.}\label{lst:m3:power-markets-design:clear}
\item \textbf{Deux zones.} La \cref{prop:m3:power-markets-design:twozones}.
\omcode{code/firm/dayahead/firm_dayahead.py}{61}{75}{Deux zones couplées par une ligne.}\label{lst:m3:power-markets-design:couple}
\item \textbf{Exécuter} \texttt{m3\_power.hour(60, 10)}, \texttt{two\_zones(3)},
\texttt{negative\_stats()} et \texttt{fig\_power.py}.
\end{tutsteps}
\textbf{Pour aller plus loin.} Relever le prix du carbone jusqu'à ce que le charbon et le gaz échangent
leur place dans l'\omterm{def:m3:power-markets-design:merit}{ordre de préséance économique}~; ajouter un \omterm{def:m3:power-markets-design:block}{ordre bloc} pour une centrale au lignite sur les
heures solaires et vérifier s'il est accepté.
\end{tutorial}

\section{Construction~: le moteur de calage day-ahead}\label{bld:m3:power-markets-design:dayahead}

\begin{build}
\bfield{Rôle} L'entreprise fictive propose des centrales et des batteries aux enchères day-ahead et
doit prévoir les prix que ses propres offres contribuent à fixer~: elle a besoin d'un moteur de calage qu'elle peut faire tourner sur ses
propres prévisions des courbes de tous les autres.
\bfield{Interface} \texttt{Order(side, price, qty)}~; \texttt{clear(orders) -> Result(price,
volume, welfare, fills)}~; \texttt{couple(zone\_a, zone\_b, capacity)}~;
\texttt{Block(name, price, qty, mtus)}~; \texttt{clear\_with\_blocks(hourly, blocks)}~;
\texttt{marginal\_cost}~; \texttt{spark\_spread}~; \texttt{dark\_spread}.
\bfield{Règles} Prix uniforme par zone et par MTU~; l'ordre marginal partiellement exécuté fixe le
prix~; des courbes qui ne se croisent pas constituent une erreur~; un bloc est tout ou rien et n'est jamais
accepté de façon paradoxale~; l'énumération se limite à de petits ensembles de blocs (l'algorithme réel
utilise la méthode par séparation et évaluation).
\bfield{Tests d'acceptation} \texttt{code/firm/dayahead/tests/}~: l'offre marginale fixe le
prix~; un prix négatif~; la congestion scinde les prix et procure la rente~; convergence avec une
ligne de grande capacité~; un bloc trop cher rejeté.
\bfield{Pour aller plus loin} Plus de deux zones sur un réseau maillé avec des contraintes fondées sur les flux~;
des MTU de 15~minutes avec des blocs horaires~; les règles de limite de variation de cours de l'encadré daté.
\end{build}

\begin{omsources}
\item EPEX SPOT et le Market Coupling Steering Committee, notices sur la MTU de 15~minutes
dans le SDAC (2025)~; ENTSO-E, Single Day-ahead Coupling.
\item ACER Decision 02/2026 on the harmonised maximum and minimum clearing price for SDAC~;
SDAC communication note du 7 mai 2026.
\item PJM, \emph{Formation of Locational Marginal Pricing and its System Energy,
Congestion and Loss Components}.
\item Supreme Court of Texas, décision de juin 2024 sur les pricing orders de la PUCT de février
2021 (d'après une note client de Reed Smith).
\item Potomac Economics (ERCOT's independent market monitor), recommandation au PUCT,
Project No.\ 51812, 4 mars 2021 (copie de l'Internet Archive).
\item Bundesnetzagentur | SMARD.de~: day-ahead prices DE-LU, generation and load, 2024--2025~; Bundesnetzagentur, communiqué de presse
sur le marché de l'électricité en 2025.
\end{omsources}

\section{Exercices}

\begin{exercise}[$\star$]\label{exo:m3:power-markets-design:1}
Offres~: 300~MWh à \qty{0}{}, 200 à \qty{40}{}, 200 à \qty{90}{}, 100 à \qty{150}{EUR/MWh}.
La demande présente des offres d'achat de 450~MWh au plafond de prix. Donner le prix et le volume accepté de chaque offre.
\end{exercise}

\begin{exercise}[$\star$]\label{exo:m3:power-markets-design:2}
Une centrale à gaz de rendement 50~\% brûle du gaz à \qty{30}{EUR/MWh} et émet 0.202~t par MWh de
gaz~; le carbone coûte \qty{70}{EUR/t}. Quel est son coût marginal~?
\end{exercise}

\begin{exercise}[$\star$]\label{exo:m3:power-markets-design:3}
Pourquoi chaque vendeur accepté reçoit-il le prix de compensation plutôt que sa propre offre~? Que
feraient les vendeurs dans un marché rémunéré au prix offert~?
\end{exercise}

\begin{exercise}[$\star\star$]\label{exo:m3:power-markets-design:4}
Dans la \cref{fig:m3:power-markets-design:zones}, quels sont le flux, les prix et la
\omterm{def:m3:power-markets-design:lmp}{rente de congestion} si la capacité de la ligne est de 3~GW~? Quelle capacité minimale
fait converger les prix~?
\end{exercise}

\begin{exercise}[$\star\star$]\label{exo:m3:power-markets-design:5}
Une ferme solaire reçoit une prime de \qty{60}{EUR/MWh} en sus du prix de marché pour chaque
MWh qu'elle produit. En dessous de quel prix s'arrête-t-elle~? À quelles heures du 11 mai 2025
se serait-elle arrêtée~?
\end{exercise}

\begin{exercise}[$\star\star$]\label{exo:m3:power-markets-design:6}
Expliquer pourquoi les prix nodaux diffèrent à l'intérieur d'une zone qui n'a qu'un seul prix zonal, et qui paie
l'écart dans un système zonal.
\end{exercise}

\begin{exercise}[$\star\star\star$]\label{exo:m3:power-markets-design:7}
\emph{Programmation.} Avec \texttt{negative\_stats}, donner le nombre d'heures négatives en 2024 et
en 2025, l'heure locale qui en compte le plus, ainsi que le prix le plus bas et son heure.
\end{exercise}

\begin{exercise}[$\star\star\star$]\label{exo:m3:power-markets-design:8}
\emph{Repérer le défaut.} «~Les prix négatifs prouvent que le marché est cassé~: personne ne devrait jamais payer
pour produire.~»
\end{exercise}

\section{Problème~: Le jour où il y a trop de soleil}

\begin{problem}[{Problème du week-end --- 11 mai 2025 pour une unité au lignite}]\label{pb:m3:power-markets-design:1}
Une unité au lignite de l'\omterm{def:m3:power-markets-design:merit}{ordre de préséance économique} de la \cref{fig:m3:power-markets-design:merit} (coût
marginal \qty{82.38}{EUR/MWh}) ne peut pas descendre en dessous de 400~MW. Le 11 mai 2025, les prix J-1 allemands
ont été négatifs de 09~h~00 à 18~h~00, heure locale.

\medskip\noindent\textbf{Partie I --- La journée.}
\begin{enumerate}
\item Pendant combien d'heures le prix a-t-il été négatif, et quelle a été sa moyenne~?
\item Quel a été le prix le plus bas, et à quel moment~?
\item À quelles heures la production solaire a-t-elle dépassé la charge~?
\item Où est allé le surplus~?
\item Pourquoi l'enchère n'a-t-elle pas été compensée à zéro~?
\end{enumerate}

\medskip\noindent\textbf{Partie II --- L'unité.}
\begin{enumerate}[resume]
\item Que perd l'unité, par rapport à une absence de production, sur une heure à charge minimale au
prix le plus bas~?
\item Que perd-elle sur les neuf heures négatives à charge minimale~?
\item Pourquoi compare-t-elle cette perte à son coût de redémarrage~?
\item Énoncer la règle de maintien en service.
\item Quelle offre d'achat exprimerait cette règle dans l'enchère~?
\end{enumerate}

\medskip\noindent\textbf{Partie III --- Les producteurs subventionnés.}
\begin{enumerate}[resume]
\item Une ferme solaire avec une prime de \qty{60}{EUR/MWh} produit-elle pendant quelles heures négatives~?
\item Quel prix un grand bloc de telles fermes met-il à l'enchère~?
\item Comment les règles de subvention ont-elles été modifiées pour limiter cet effet~?
\item Pourquoi le stockage par batteries fait-il monter les prix de midi~?
\item Quels producteurs tirent profit des prix négatifs~?
\end{enumerate}

\medskip\noindent\textbf{Partie IV --- Jugement.}
\begin{enumerate}[resume]
\item Qu'aurait fait une interconnexion deux fois plus grande au prix allemand ce jour-là~?
\item Pourquoi les prix négatifs se concentrent-ils au printemps~?
\item Quelle est la valeur d'un consommateur flexible un tel jour~?
\item Énoncer le \emph{résultat nommé}~: le coût de redémarrage en dessous duquel l'unité doit s'arrêter.
\item En une phrase~: que signale un prix négatif de l'électricité~?
\end{enumerate}
\end{problem}

\section{Questions d'entretien}

\begin{interviewq}[$\star$ \iqroles{trader}]\label{iq:m3:power-markets-design:1}
Comment le prix J-1 de l'électricité d'une zone européenne est-il fixé~?
\end{interviewq}

\begin{interviewq}[$\star$ \iqroles{trader, chercheur}]\label{iq:m3:power-markets-design:2}
Qu'est-ce que l'\omterm{def:m3:power-markets-design:merit}{ordre de préséance économique}, et comment l'ajout d'éolien modifie-t-il les prix~?
\end{interviewq}

\begin{interviewq}[$\star\star$ \iqroles{chercheur}]\label{iq:m3:power-markets-design:3}
Vous voulez prévoir le prix J-1 de demain pour chaque heure. Quelles entrées comptent, et
laquelle est la plus importante au printemps~?
\end{interviewq}

\begin{interviewq}[$\star\star$ \iqroles{trader, risque}]\label{iq:m3:power-markets-design:4}
Pourquoi un prix de l'électricité peut-il être négatif, et jusqu'où peut-il descendre~?
\end{interviewq}

\begin{interviewq}[$\star\star$ \iqroles{chercheur}]\label{iq:m3:power-markets-design:5}
Comparez la tarification zonale et la tarification nodale du point de vue du propriétaire d'une ferme éolienne.
\end{interviewq}

\begin{interviewq}[$\star\star\star$ \iqroles{développeur, chercheur}]\label{iq:m3:power-markets-design:6}
Concevez un simulateur de compensation J-1 pour dix zones et 96 MTU avec des \omterm{def:m3:power-markets-design:block}{ordres bloc}. Comment
le garder assez rapide pour exécuter des milliers de scénarios~?
\end{interviewq}
